\documentclass[11pt]{amsart}
\usepackage{amssymb,amsmath,amsthm,mathtools}
\usepackage[margin=1.1in]{geometry}
\usepackage{enumitem}
\usepackage{booktabs}
\usepackage{array}
\usepackage[hidelinks]{hyperref}

\newtheorem{theorem}{Theorem}[section]
\newtheorem{proposition}[theorem]{Proposition}
\newtheorem{lemma}[theorem]{Lemma}
\newtheorem{corollary}[theorem]{Corollary}
\newtheorem{conjecture}[theorem]{Conjecture}
\theoremstyle{definition}
\newtheorem{definition}[theorem]{Definition}
\newtheorem{question}[theorem]{Question}
\theoremstyle{remark}
\newtheorem{remark}[theorem]{Remark}

\newcommand{\QQ}{\mathbb{Q}}
\newcommand{\RR}{\mathbb{R}}
\newcommand{\ZZ}{\mathbb{Z}}
\newcommand{\CC}{\mathbb{C}}
\newcommand{\PP}{\mathbb{P}}
\newcommand{\cO}{\mathcal{O}}
\newcommand{\Hdg}{\operatorname{Hdg}}
\newcommand{\Alg}{\operatorname{Alg}}
\newcommand{\End}{\operatorname{End}}
\newcommand{\Hg}{\operatorname{Hg}}
\newcommand{\SU}{\operatorname{SU}}
\newcommand{\U}{\operatorname{U}}
\newcommand{\Pic}{\operatorname{Pic}}
\newcommand{\cl}{\operatorname{cl}}
\newcommand{\ch}{\operatorname{ch}}
\newcommand{\Sh}{\mathrm{Sh}}
\newcommand{\HC}{\mathrm{HC}}

\title[The Hodge conjecture: an attempt in full]{The Hodge conjecture: an
attempt in full,\\ and the statements at which it stops}
\author{Deep Bhattacharjee}
\date{October 2026}

\begin{document}

\begin{abstract}
We follow every route to the Hodge conjecture that the present literature
makes available, prove each reduction that can be proved, and stop at the
first statement on each route that cannot. The Hodge conjecture is not
proved here. Every Hodge class of degree two, of codegree two, or on a
variety of dimension at most three is algebraic, and the conjecture descends
along surjective morphisms. By the author's earlier work it is equivalent to
two statements: the conjecture for abelian varieties, and the conjecture
modulo abelian varieties. For abelian varieties of Weil type the problem comes down to one
class on a general member of a family, and we show that the cycles Schoen
constructed on Prym varieties of triple covers give, on each such Prym
variety, an explicit subvariety of half dimension whose class is a positive
multiple of the polarization power plus a nonzero Weil class. If that
subvariety were semiregular, the Weil classes would be algebraic on every
member of the split Weil families of $\QQ(\sqrt{-3})$ in every dimension. We
cannot decide whether it is. We also show that every abelian variety of Weil
type for $\QQ(\sqrt{-3})$ is, up to isogeny, a factor of such a Prym
variety, so that its Weil classes are algebraic exactly when those of the
complementary factor are; a count of dimensions indicates that this cannot
reach the general member in dimension eight or more. For the remaining
classes on abelian varieties,
and for varieties that are not of abelian type, we record what is proved and
locate the first open case of each: the exceptional classes on the square of
a Mumford fourfold, and the square of a K3 surface with real
multiplication.
\end{abstract}

\maketitle

\section{Introduction}

Let $X$ be a smooth complex projective variety of dimension $n$. A class in
$H^{2p}(X,\QQ)$ is a \emph{Hodge class} if its image in $H^{2p}(X,\CC)$ lies
in $H^{p,p}(X)$, and it is \emph{algebraic} if it is a rational combination of
classes of subvarieties of codimension $p$. Algebraic classes are Hodge
classes. The Hodge conjecture says that the converse holds.

\begin{conjecture}[Hodge]
Every Hodge class on a smooth complex projective variety is algebraic.
\end{conjecture}

This paper is an attempt at the conjecture in full. By that we mean
something specific. We take every route that the literature makes
available, push it as far as it can be pushed with complete proofs, and
name the first statement on it that we cannot prove. We do not prove the
conjecture, and nothing here should be read as a claim that it is close to
proved. What the attempt produces is a map: a short list of statements, each
open, each precise, from which the conjecture would follow, together with
an account of why the obvious ways of proving each of them fail.

Section~\ref{sec:reductions} proves the reductions that hold for every
variety: degree two, codegree two, dimension at most three, descent along
surjective morphisms, and the reduction of every degree to the middle one
under the Lefschetz standard conjecture. Section~\ref{sec:shape} recalls
from \cite{BhH8} that, beyond these, the conjecture is equivalent to two
statements, the conjecture for abelian varieties and the conjecture modulo
abelian varieties, and lists the other minimal sets of hypotheses from which
it follows. Sections~\ref{sec:abelian} and~\ref{sec:weil} attack abelian
varieties. Section~\ref{sec:weil} contains the main new theorem of the
paper, Theorem~\ref{thm:schoensub}, which turns Schoen's cycles into a
single explicit subvariety of the right shape for a deformation argument,
and isolates the question on which the argument turns.
Section~\ref{sec:beyondweil} treats the classes on abelian varieties that
divisors and Weil classes do not generate, and Section~\ref{sec:nonabelian}
the varieties that are not of abelian type. Section~\ref{sec:spread} tries
to spread Schoen's construction beyond the Prym varieties. It proves that
every abelian variety of Weil type for $\QQ(\sqrt{-3})$ is, up to isogeny, a
factor of the Prym variety of an \'etale triple cover, so that its Weil
classes are algebraic exactly when those of the complementary factor are
(Theorem~\ref{thm:transfer}); it shows by a count why this cannot by itself
reach the general member once the dimension is at least eight; and it
explains why the construction, carried to other varieties, leads back to
(F2) and (F3$'$). Section~\ref{sec:end} collects what is proved and what is
not.

Table~\ref{tab:routes} summarises where each route stops.

\begin{table}[ht]
\centering
\small
\begin{tabular}{@{}>{\raggedright\arraybackslash}p{4.4cm}>{\raggedright\arraybackslash}p{4.6cm}>{\raggedright\arraybackslash}p{4.6cm}@{}}
\toprule
Route & What is proved & Where it stops\\
\midrule
Divisors and hard Lefschetz & degrees $2$ and $2n-2$; $\dim X\le3$ & degree $2p$ with $2\le p\le n-2$\\
Hyperplane sections & reduction to middle degree, given the Lefschetz standard conjecture & the Lefschetz standard conjecture\\
Weil classes, deformation & the algebraic locus is dense and all or meagre; one semiregular cycle closes it & no semiregular cycle is known; Schoen's subvariety is undecided\\
Weil classes, degree bounds & equivalent to a uniform degree bound on a dense set & no such bound is known\\
Weil classes, Prym factors & for $\QQ(\sqrt{-3})$, algebraic on $A$ if and only if on a complement $A'$ & no general complement is known; a count says none exists for $\dim A\ge8$\\
Classes beyond Weil & algebraic at CM points of Mumford families & one Kuga--Satake class\\
Not of abelian type & every variety reached from abelian varieties; many hyper-K\"ahler varieties & the square of a K3 surface with real multiplication\\
\bottomrule
\end{tabular}
\caption{The routes and their stopping points.}\label{tab:routes}
\end{table}

\subsection*{Conventions}
Varieties are smooth, complex, projective and connected unless said
otherwise. For such an $X$ we write $\Hdg^{p}(X)\subseteq H^{2p}(X,\QQ)$ for
the Hodge classes and $\Alg^{p}(X)\subseteq\Hdg^{p}(X)$ for the algebraic
ones, and $\HC(X,p)$ for the statement $\Alg^{p}(X)=\Hdg^{p}(X)$. We write
$\HC(X)$ when this holds for every $p$.

\section{Reductions that hold for every variety}\label{sec:reductions}

\subsection{Degree two}

\begin{theorem}[Lefschetz {\cite{Lef24}}, Kodaira--Spencer {\cite{KS53}}]\label{thm:lef11}
$\HC(X,1)$ holds for every $X$.
\end{theorem}

\begin{proof}
The exponential sequence $0\to\ZZ\to\cO_{X}\to\cO_{X}^{\times}\to0$ gives an
exact sequence
\[
  H^{1}(X,\cO_{X}^{\times})\xrightarrow{\;c_{1}\;}H^{2}(X,\ZZ)\longrightarrow
  H^{2}(X,\cO_{X}).
\]
The second map is the composite of $H^{2}(X,\ZZ)\to H^{2}(X,\CC)$ with the
projection onto $H^{0,2}(X)$ given by the Hodge decomposition. An integral
class of type $(1,1)$ has zero projection to $H^{0,2}$, so it is the first
Chern class of a holomorphic line bundle. By the GAGA theorem of Serre that
line bundle is algebraic, and its first Chern class is the class of a divisor.
A rational Hodge class has an integral multiple, which is therefore the class
of a divisor.
\end{proof}

\subsection{Hard Lefschetz}

Let $h\in H^{2}(X,\QQ)$ be the class of a hyperplane section, and $L$ the
operator of cup product with $h$.

\begin{proposition}\label{prop:hardlef}
For $p<n/2$, $\HC(X,p)$ implies $\HC(X,n-p)$.
\end{proposition}

\begin{proof}
By the hard Lefschetz theorem $L^{n-2p}\colon H^{2p}(X,\QQ)\to
H^{2n-2p}(X,\QQ)$ is an isomorphism. It maps $H^{p,p}$ to $H^{n-p,n-p}$, so
it maps $\Hdg^{p}(X)$ onto $\Hdg^{n-p}(X)$: if $\beta$ is a Hodge class of
degree $2n-2p$ and $\beta=L^{n-2p}\alpha$, then $L^{n-2p}$ maps the $(r,s)$
components of $\alpha$ to the $(r+n-2p,s+n-2p)$ components of $\beta$, so
$\alpha$ has type $(p,p)$. Cup product with $h$ sends the class of a
subvariety $Z$ to the class of $Z\cap H$ for a general hyperplane $H$, so
$L^{n-2p}$ maps $\Alg^{p}(X)$ into $\Alg^{n-p}(X)$. Hence
$\Hdg^{n-p}(X)=L^{n-2p}\Hdg^{p}(X)=L^{n-2p}\Alg^{p}(X)\subseteq\Alg^{n-p}(X)$.
\end{proof}

\begin{corollary}\label{cor:smalldim}
$\HC(X,p)$ holds for $p\in\{0,1,n-1,n\}$. In particular $\HC(X)$ holds when
$\dim X\le3$.
\end{corollary}

\begin{proof}
For $p=0$ and $p=n$ the Hodge classes are multiples of the fundamental class
and of the class of a point. Theorem~\ref{thm:lef11} gives $p=1$, and
Proposition~\ref{prop:hardlef} with $p=1$ gives $p=n-1$. When $n\le3$ every
$p$ is one of these.
\end{proof}

\subsection{Descent}

\begin{proposition}\label{prop:descent}
Let $f\colon Y\to X$ be a surjective morphism of smooth projective
varieties, and $r=\dim Y-\dim X$. Then $\HC(Y,p+r)$ implies $\HC(X,p)$. In
particular $\HC(Y)$ implies $\HC(X)$.
\end{proposition}

\begin{proof}
Let $h_{Y}$ be the class of a hyperplane section of $Y$. The push-forward
$f_{*}(h_{Y}^{r})\in H^{0}(X,\QQ)$ is the degree $d$ of a general fibre
of $f$ against $h_{Y}^{r}$, which is a positive integer. For
$\alpha\in\Hdg^{p}(X)$ the class $f^{*}\alpha\cup h_{Y}^{r}$ is a Hodge class
of degree $2(p+r)$ on $Y$, hence algebraic, and by the projection formula
$f_{*}(f^{*}\alpha\cup h_{Y}^{r})=\alpha\cup f_{*}(h_{Y}^{r})=d\,\alpha$.
Push-forward of cycles is compatible with the cycle class, so $d\alpha$, and
hence $\alpha$, is algebraic.
\end{proof}

Taking $Y=X\times\PP^{1}$ shows that the conjecture for $X\times\PP^{1}$ gives
it for $X$. This is the reason, explained in \cite{BhH8}, that the
conjecture for the varieties that are not abelian already contains the
conjecture for abelian varieties.

\subsection{The reduction to middle degree}

The remaining degrees are $2\le p\le n-2$, and by
Proposition~\ref{prop:hardlef} it is enough to treat $2\le p<n/2$ and, when
$n$ is even, $p=n/2$. The natural way to move from degree $2p$ on $X$ to the
middle degree of a smaller variety goes through hyperplane sections, and it
needs one more statement. Recall that the \emph{Lefschetz standard
conjecture} $B(X)$ asks that the inverse of $L^{n-2p}$, defined on
$H^{2n-2p}(X,\QQ)$, be induced by an algebraic cycle on $X\times X$ for every
$p$.

\begin{proposition}\label{prop:middle}
Let $2p<n$ and let $i\colon Y\hookrightarrow X$ be a smooth complete
intersection of $n-2p$ hyperplane sections. If $B(X)$ and $\HC(Y,p)$ hold,
then $\HC(X,p)$ holds.
\end{proposition}

\begin{proof}
Let $\alpha\in\Hdg^{p}(X)$. Then $i^{*}\alpha\in\Hdg^{p}(Y)$ is algebraic, say
$i^{*}\alpha=\cl(Z)$, and $i_{*}i^{*}\alpha=\alpha\cup[Y]=L^{n-2p}\alpha$, so
$L^{n-2p}\alpha=\cl(i_{*}Z)$ is algebraic. The inverse of $L^{n-2p}$ is
induced by an algebraic correspondence, which maps algebraic classes to
algebraic classes, so $\alpha$ is algebraic.
\end{proof}

The Lefschetz standard conjecture is itself open. It is known for curves and
surfaces, for abelian varieties by Lieberman \cite{Lie68,Kle68}, and for
hyper-K\"ahler varieties of $K3^{[n]}$ type by Charles and Markman
\cite{CM13}. It follows from the Hodge conjecture, so
Proposition~\ref{prop:middle} is a reduction and not a proof. Without it we
do not know how to get from degree $2p$ on $X$ to any smaller variety, and
this is the first place where the general argument stops.

\subsection{What any proof must use}

Two counterexamples constrain the shape of a proof. The integral version of
the conjecture fails: Atiyah and Hirzebruch found torsion classes that are
not algebraic \cite{AH62}, and Koll\'ar found non-torsion integral Hodge
classes with no algebraic representative although a multiple has one
\cite{Kol92}. So a proof has to use rational coefficients in an essential
way. The version for compact K\"ahler manifolds fails too: Voisin found
complex tori with Hodge classes that are not combinations of Chern classes of
coherent sheaves \cite{Voi02b}. So a proof has to use projectivity, and an
argument that would work on every compact K\"ahler manifold cannot be right.
Both facts are used below: Voisin's tori are what rule out one form of the
deformation argument in Section~\ref{sec:weil}.

\section{The shape of what remains}\label{sec:shape}

The following two statements were isolated in \cite[Section~14]{BhH8}.

\begin{definition}\label{def:f2f3}
\begin{enumerate}[label=(\roman*)]
\item (F2) For every abelian variety $A$, the Hodge classes of $A$ that lie
outside the subring generated by divisor classes and the Weil classes of
quotients of $A$ are algebraic.
\item (F3$'$) For every $X$, every Hodge class on $X$ is algebraic modulo
the images of Hodge classes of abelian varieties under algebraic
correspondences.
\end{enumerate}
\end{definition}

\begin{theorem}[{\cite[Theorem 14.1(iv), Propositions on (F2) and (F3$'$)]{BhH8}}]\label{thm:shape}
\begin{enumerate}[label=(\alph*)]
\item The Hodge conjecture is equivalent to the conjunction of
\textup{(F2)} and \textup{(F3$'$)}.
\item \textup{(F2)} is equivalent to the Hodge conjecture for abelian
varieties.
\item Over the implications proved in \cite{BhH8} and in the literature,
the minimal sets of statements that imply the conjecture are
$\{\mathrm{F3}\}$, $\{L,M\}$, $\{L,\mathrm{F3}'\}$, $\{V,\mathrm{F3}'\}$ and
$\{\mathrm{F2},\mathrm{F3}'\}$, where \textup{(F3)} is the conjecture for
varieties that are not abelian, $L$ is the Lefschetz standard conjecture for
every variety, $M$ says that every Hodge class is motivated in the sense of
Andr\'e \cite{And96}, and $V$ is the variational Hodge conjecture for
algebraic classes.
\end{enumerate}
\end{theorem}

Part~(a) has a simple meaning. Given (F2), every Hodge class on an abelian
variety is algebraic, and an algebraic correspondence carries algebraic
classes to algebraic classes; so (F3$'$) then says that every Hodge class on
every variety is algebraic. The attempt therefore splits into two halves:
the conjecture for abelian varieties, which is (F2) and is the subject of
Sections~\ref{sec:abelian} to~\ref{sec:beyondweil}, and the conjecture modulo
abelian varieties, which is (F3$'$) and is the subject of
Section~\ref{sec:nonabelian}. We note that (F3) alone is the whole
conjecture, by Proposition~\ref{prop:descent} applied to $A\times\PP^{1}\to A$,
so the split is a real division of the problem and not a reformulation of
it.

\section{Abelian varieties}\label{sec:abelian}

Let $A$ be an abelian variety of dimension $g$ and $V=H^{1}(A,\QQ)$. Then
$H^{*}(A,\QQ)=\bigwedge^{*}V^{\vee}$, and the Hodge structure on $V$ is given by
a homomorphism $h\colon\U(1)\to\operatorname{GL}(V_{\RR})$. The \emph{Hodge
group} $\Hg(A)$ is the smallest algebraic subgroup of $\operatorname{GL}(V)$
defined over $\QQ$ whose real points contain the image of $h$.

\begin{proposition}[{\cite[Section~4]{BhH8}, \cite{Del82}}]\label{prop:invariants}
The Hodge classes of $A$ in every degree are exactly the elements of
$\bigwedge^{*}V^{\vee}$ fixed by $\Hg(A)$.
\end{proposition}

The \emph{Lefschetz group} of $A$ is the centraliser of $\End^{0}(A)$ in the
symplectic group of a polarization. Its invariants are generated by divisor
classes, and the Hodge group is contained in it. When the two are equal,
every Hodge class is a polynomial in divisor classes, and is algebraic. The
Hodge classes that are not generated by divisors exist exactly when the Hodge
group is strictly smaller than the Lefschetz group, and they come in two
kinds.

The first kind are Weil's classes \cite{Wei77}. Let $K$ be an imaginary
quadratic field acting on $A$ of dimension $2n$ so that both embeddings of
$K$ occur with multiplicity $n$ on $H^{1,0}(A)$. Then
$W=\bigwedge^{2n}_{K}V^{\vee}$ is a $\QQ$-plane of Hodge classes in
$H^{2n}(A,\QQ)$, the \emph{Weil plane}, and for a general such $A$ it is not
generated by divisors. These are the subject of Section~\ref{sec:weil}.

The second kind appear on abelian varieties whose Hodge group is smaller than
the Lefschetz group for a reason other than a $K$-action. The simplest are
Mumford's fourfolds \cite{Mum69}: abelian fourfolds with
$\End(A)=\ZZ$ whose Hodge group is a form of $\mathrm{SL}_{2}^{3}$ acting on
$V$ through the tensor product of the three standard representations. They
have no exceptional Hodge classes on $A$ itself, but $A\times A$ has two.
These are the subject of Section~\ref{sec:beyondweil}.

The known cases of the conjecture for abelian varieties are those where the
Hodge ring is generated by divisors (for instance when $\dim A\le3$, or when
$A$ is a power of an elliptic curve without complex multiplication), the
fourfolds and some sixfolds of Weil type listed in Section~\ref{sec:weil},
and the products and powers covered by \cite{MZ99}, \cite{RM08} and
\cite{Fl25}.

\section{Weil classes}\label{sec:weil}

\subsection{The reduction to one class}

Fix an imaginary quadratic field $K$, an integer $n\ge2$ and a nondegenerate
$K$-hermitian form of signature $(n,n)$, with discriminant $\delta$. The
abelian varieties of Weil type with these data form a family $\Sh$ over a
connected base whose universal cover is the bounded symmetric domain of
$\SU(n,n)$, of dimension $n^{2}$. A very general member has Hodge group
$\SU(V,H)$, and for such a member every Hodge class is a polynomial in the
polarization class $\eta$ except in degree $2n$, where
\[
  \Hdg^{n}(A)=\QQ\,\eta^{n}\oplus W
\]
\cite{BhH8}, \cite[Theorem~4]{Wei77}. Multiplication by a suitable
element of $K$ moves one nonzero Weil class onto a second one that spans the
Weil plane with it, so on such a member the whole conjecture is the
algebraicity of a single Weil class.

\subsection{The algebraic locus}

Let $\Sigma\subseteq\Sh$ be the set of members on which the Weil class is
algebraic. The following is proved in \cite{BhH8} and in the companion note
\cite{BhW}; we list it to fix what the deformation route already has.

\begin{theorem}[{\cite{BhH8,BhW}}]\label{thm:locus}
\begin{enumerate}[label=(\alph*)]
\item $\Sigma$ is a countable union of closed algebraic subsets of $\Sh$,
it is dense, and it is either all of $\Sh$ or meagre.
\item $\Sigma=\Sh$ if and only if the Weil class has algebraic
representatives of bounded degree on a Zariski-dense set of members.
\item If at one member there is a local complete intersection, or a perfect
complex, whose class (respectively Chern character) in degree $2n$ is a
nonzero Weil class plus a multiple of $\eta^{n}$, with higher terms in the
span of the powers of $\eta$, and whose semiregularity map is injective,
then $\Sigma=\Sh$.
\item No object whose Chern character in degree $2n$ is a pure Weil class is
semiregular.
\item At a member with Hodge group $\SU(V,H)$, no cycle or complex built from
divisors and homomorphisms by the usual operations represents a Weil class,
and neither does any tautological cycle of a curve with an automorphism of
order three.
\end{enumerate}
\end{theorem}

Part~(d) is where the K\"ahler counterexample enters: a semiregular object
would deform with the class, and a pure Weil class stays of Hodge type on a
family of complex tori that contains Voisin's tori, which carry no
subvarieties and no coherent sheaves with the required Chern character. Part
(c) does not meet this obstruction, because $\eta^{n}$ stops being of Hodge
type as soon as the torus leaves the polarized family. So the deformation
route asks for a semiregular object at some member whose class is a nonzero
Weil class plus a multiple of $\eta^{n}$. Part~(e) says that such an object
cannot be built from the obvious ingredients at a member with full Hodge
group. At special members it can be, but transporting it to the general
member by Hecke correspondences makes its degree grow without bound, and by
part~(b) bounded degree is exactly what is needed.

\subsection{The known cases}

The Weil class is known to be algebraic on every member in the following
families: Weil fourfolds for $K=\QQ(\sqrt{-3})$ of every discriminant
\cite{Sch88,Sch98}; a family of Weil sixfolds for $\QQ(\sqrt{-3})$
\cite{Sch98} and one for $\QQ(\sqrt{-1})$ \cite{Koi04}, both obtained from
Prym varieties of cyclic covers of curves of genus four; Weil fourfolds of
discriminant one for every $K$, and their powers \cite{Fl25}. Markman's
preprints extend this to every Weil fourfold and every split Weil sixfold
\cite{Mar25a}. For $2n\ge8$ no family is known in which the Weil class is
algebraic on every member.

\subsection{Schoen's cycles as one subvariety}

The only cycles known to represent Weil classes at members with full Hodge
group in dimension $2n\ge8$ are Schoen's. We recall his construction and
then extract from it a single subvariety of the shape that
Theorem~\ref{thm:locus}(c) asks for.

Let $X$ be a curve of genus $g=n+1\ge3$, let $L\in\Pic^{0}(X)$ have exact
order three, and let $\pi\colon C\to X$ be the \'etale cyclic triple cover
defined by $L$, with covering automorphism $\sigma$. The curve $C$ has genus
$3g-2$. In the Jacobian $J=J(C)$ let $N=1+\sigma+\sigma^{2}$ and
$P=3-N\in\End(J)$, and let $B=P(J)$, the \emph{Prym variety} of the cover.
Then $N^{2}=3N$, so $P$ is multiplication by $3$ on $B$, and $\sigma$ acts on
$B$ with minimal polynomial $x^{2}+x+1$, which makes $B$ an abelian variety
of dimension $2n$ with an action of $K=\QQ(\zeta)$, $\zeta=e^{2\pi i/3}$. Its
hermitian form is split \cite{Loo97}, and both embeddings of $K$ occur with
multiplicity $n$ on $H^{1,0}(B)$ \cite[Lemma~1.6a]{Sch88}, so $B$ is of Weil
type. For very general $(X,L)$ its Hodge group is $\SU(V,H)$, by Looijenga's
computation of the monodromy \cite{Loo97} and Andr\'e's theorem
\cite{And92}.

Put $h=2g-2=2n$, fix $c_{0}\in C$, and let
\[
  \Psi\colon C^{h}\to B,\qquad (c_{1},\dots,c_{h})\mapsto
  P\Bigl(\,\sum_{i}(c_{i}-c_{0})\Bigr).
\]
Let $\Lambda\subseteq C^{h}$ be the set of $h$-tuples with
$\pi(c_{1})+\dots+\pi(c_{h})\in|K_{X}|$. Since $|K_{X}|\cong\PP^{n}$ and the
map $C^{h}\to X^{(h)}$ is finite, $\Lambda$ has pure dimension $n$.

Schoen's construction \cite[Section~2, case $r=0$]{Sch88} runs as follows. Let
$G'\subseteq\operatorname{Aut}(C^{h})$ be generated by the permutations of the
factors and by the maps $\sigma^{v_{1}}\times\dots\times\sigma^{v_{h}}$ with
$\sum v_{i}\equiv0\pmod3$, and let $W_{0}=C^{h}/G'$ with quotient map $\rho$.
The map $W_{0}\to X^{(h)}$ is an \'etale cover of degree three, and over
$|K_{X}|$, which is simply connected, it splits into three components $Q_{0}$,
$Q_{1}$, $Q_{2}$, permuted by $\delta=\sigma\times1\times\dots\times1$. For
each nontrivial character $\chi$ of $\ZZ/3$ the cycle
$z_{\chi}=\sum_{t}\chi(-t)\,\delta^{t}_{*}Q_{0}$ has nonzero class
\cite[proof of Theorem~2.0]{Sch88}. Patel and Zhang reach the same three
components by base change of the Abel--Jacobi map along an isogeny onto
$J(X)$ \cite[Theorem~1.1]{PZ25}. Let $\Lambda_{0}=\rho^{-1}(Q_{0})\subseteq
\Lambda$.

\begin{theorem}\label{thm:schoensub}
In this situation:
\begin{enumerate}[label=(\alph*)]
\item the class $\Psi_{*}[\Lambda_{0}]\in H^{2n}(B,\QQ)$ pairs nontrivially
with the Weil plane of $B$;
\item there is an irreducible component $\Lambda'$ of $\Lambda_{0}$ that
maps generically finitely onto an $n$-dimensional subvariety $Y=\Psi(\Lambda')$
of $B$ whose class pairs nontrivially with the Weil plane;
\item if $\Hg(B)=\SU(V,H)$, which holds for very general $(X,L)$, then
\[
  \cl(Y)=c\,\eta^{n}+w,\qquad c\in\QQ_{>0},\quad 0\ne w\in W.
\]
\end{enumerate}
\end{theorem}

\begin{proof}
We work with complex coefficients and write $\bar\chi$ for the conjugate
character. Write $H^{1}(B,\CC)=V_{\chi}\oplus V_{\bar\chi}$ for the
eigenspaces of $\sigma^{*}$, each of dimension $h$. The Weil plane, tensored
with $\CC$, is $\bigwedge^{h}V_{\chi}\oplus\bigwedge^{h}V_{\bar\chi}$.

\emph{Step 1: the pull-back of the Weil plane.} Write $\Psi=P'\circ\Xi$, where
$\Xi\colon C^{h}\to J$ is $(c_{i})\mapsto\sum(c_{i}-c_{0})$ and $P'\colon J\to
B$ is $P$ with its target restricted. If $\iota\colon B\to J$ is the
inclusion, then $P'\iota$ is multiplication by $3$ on $B$, so $P'^{*}$ is
injective on $H^{1}$, and it carries $V_{\chi}$ into the $\chi$-eigenspace of
$H^{1}(C,\CC)$, which $\Xi^{*}$ identifies with a summand of
$H^{1}(C^{h},\CC)$ on each factor. For a basis $y_{1},\dots,y_{h}$ of
$V_{\chi}$ let $\tilde y_{j}\in H^{1}(C,\CC)_{\chi}$ be the corresponding
classes on $C$. Then
\[
  u_{\chi}:=\Psi^{*}(y_{1}\wedge\dots\wedge y_{h})
  =\prod_{j=1}^{h}\Bigl(\sum_{i=1}^{h}p_{i}^{*}\tilde y_{j}\Bigr),
\]
where $p_{i}$ is the $i$-th projection. A product $\tilde y_{j}\cup\tilde
y_{k}$ lies in the $\chi^{2}$-eigenspace of $H^{2}(C,\CC)$. Since $\sigma$
acts trivially on $H^{2}(C,\CC)$ and $\chi^{2}\ne1$, that eigenspace is zero,
so the product vanishes. Expanding the product, only the terms with one
factor $p_{i}^{*}\tilde y_{j}$ for each $i$ survive, and
$u_{\chi}=\sum_{s}\pm\prod_{i}p_{i}^{*}\tilde y_{s(i)}$, the sum over
permutations $s$. By the K\"unneth formula $u_{\chi}\ne0$.

From here on automorphisms act on cohomology by pull-back. The class
$u_{\chi}$ is invariant under $G'$: the permutations fix it because $\Psi$ is
symmetric, and each $\sigma^{v_{1}}\times\dots\times\sigma^{v_{h}}$ multiplies
it by $\chi(\sum v_{i})=1$. The map $\delta$ normalizes $G'$, and $\delta^{*}$
multiplies $u_{\chi}$ by $\chi(1)$. Let $U_{\chi}$ be the part of
$H^{h}(C^{h},\CC)^{G'}$ on which $\delta^{*}$ acts by $\chi(1)$. It is
one-dimensional \cite[Lemma~1.2]{Sch88}. (Directly: a K\"unneth product of
eigenclasses is fixed by every $\sigma^{v_{1}}\times\dots\times\sigma^{v_{h}}$
with $\sum v_{i}\equiv0$ only if all its factors lie in one eigenspace of
$\sigma^{*}$, and $\delta^{*}$ then acts by $\chi(1)$ only if that eigenspace
is $H^{1}(C,\CC)_{\chi}$, which has dimension $2g-2=h$; the symmetric part is
its top exterior power.) So $u_{\chi}$ spans $U_{\chi}$. In the same way
$u_{\bar\chi}=\Psi^{*}(\bar y_{1}\wedge\dots\wedge\bar y_{h})$ spans
$U_{\bar\chi}$.

\emph{Step 2: $\Psi$ is generically finite.} At a point
$(c_{1},\dots,c_{h})$ the differential of $\Psi$ is dual to the map that
sends a holomorphic $1$-form $\omega$ on $B$, viewed as a form on $C$ in the
nontrivial eigenspaces, to $(\omega(c_{1}),\dots,\omega(c_{h}))$. These forms
make up a space of dimension $2n=h$, and $h$ general points impose independent
conditions on it, so the differential is an isomorphism at a general point.
Hence $\Psi$ is surjective and generically finite, of some degree $e>0$, and
\[
  \int_{C^{h}}u_{\chi}\cup u_{\bar\chi}
  =e\int_{B}(y_{1}\wedge\dots\wedge y_{h})\wedge(\bar y_{1}\wedge\dots\wedge\bar y_{h})\ne0.
\]

\emph{Step 3: the pairing.} The class $\lambda_{0}=[\Lambda_{0}]=\rho^{*}[Q_{0}]$
is $G'$-invariant; $\rho$ is unramified at the general point of
$\Lambda_{0}$, since $G'$ acts freely on tuples whose images in $X$ are
distinct. Since $\delta^{*}=(\delta^{-1})_{*}$, the component of
$\lambda_{0}$ on which $\delta^{*}$ acts by $\chi(1)$ is
\[
  \tfrac13\sum_{t}\chi(t)\,\delta^{t}_{*}\lambda_{0}
  =\tfrac13\,\rho^{*}[z_{\bar\chi}],
\]
which is nonzero by Schoen's theorem applied to $\bar\chi$, because
$\rho^{*}$ is injective in rational cohomology. It lies in $U_{\chi}$, so it
equals $a\,u_{\chi}$ with $a\ne0$. Because $\delta$ preserves the
intersection pairing, a class on which $\delta^{*}$ acts by $\psi(1)$ pairs
to zero with $u_{\bar\chi}$ unless $\psi=\chi$. Hence
\[
  \int_{C^{h}}\lambda_{0}\cup u_{\bar\chi}=a\int_{C^{h}}u_{\chi}\cup u_{\bar\chi}\ne0,
\]
and by the projection formula
$\int_{B}\Psi_{*}\lambda_{0}\cup(\bar y_{1}\wedge\dots\wedge\bar y_{h})\ne0$.
This is~(a).

\emph{Step 4: one component.} The class $\Psi_{*}\lambda_{0}$ is
$\sum_{k}e_{k}\,\cl(\Psi(\Lambda_{k}))$, the sum over the irreducible
components $\Lambda_{k}$ of $\Lambda_{0}$ on which $\Psi$ is generically
finite, with $e_{k}$ the degree of $\Lambda_{k}\to\Psi(\Lambda_{k})$; the
other components contribute nothing. By~(a) one term pairs nontrivially with
the Weil plane. This is~(b).

\emph{Step 5: the shape of the class.} If $\Hg(B)=\SU(V,H)$ then
$\cl(Y)\in\Hdg^{n}(B)=\QQ\eta^{n}\oplus W$, so $\cl(Y)=c\eta^{n}+w$. The
centre $\U(1)\subseteq\U(V,H)$, acting by scalars of $K\otimes\RR=\CC$ of
absolute value one, fixes $\eta$ and the orientation class, and acts on the
two lines of $W\otimes\CC$ by $\lambda^{\pm2n}$. So $\int\eta^{n}\cup w'=0$
for every $w'\in W$. By~(b), $w\ne0$. Finally
$\int\cl(Y)\cup\eta^{n}$ is the degree of $Y$ with respect to the ample class
$\eta$, which is positive, and it equals $c\int\eta^{2n}$, so $c>0$.
\end{proof}

\begin{corollary}\label{cor:schoensemireg}
Let $n\ge4$. Suppose that for one very general $(X,L)$ of genus $n+1$ the
subvariety $Y$ of Theorem~\ref{thm:schoensub}, or its structure sheaf as a
perfect complex on $B$, is semiregular. Then the Weil classes are algebraic on
every member of the split Weil family of $\QQ(\sqrt{-3})$ in dimension $2n$.
\end{corollary}

\begin{proof}
If $Y$ is a local complete intersection, its class is
$c\eta^{n}+w$ with $w\ne0$, which is the shape required by
Theorem~\ref{thm:locus}(c). For the structure sheaf, $\ch(\cO_{Y})$ has
$\cl(Y)$ in degree $2n$ and Hodge classes in higher degrees, which at a
member with Hodge group $\SU(V,H)$ are multiples of powers of $\eta$ by the
description of the Hodge ring. In either case Theorem~\ref{thm:locus}(c)
gives $\Sigma=\Sh$.
\end{proof}

\subsection{Where the Weil route stops}

Everything now rests on one question.

\begin{question}\label{q:schoen}
For $n\ge4$ and very general $(X,L)$, is the subvariety $Y\subseteq B$ of
Theorem~\ref{thm:schoensub} semiregular? Equivalently, at first order:
does $Y$ deform along every one of the $n^{2}$ directions of the Weil family
at $B$, and not only along the $3n$ directions that come from deforming
$(X,L)$?
\end{question}

We have not been able to decide it, and the closest analogue points the wrong
way. The Abel--Jacobi curve $C\subseteq J(C)$ is built from the same map as
$Y$. Its class $\theta^{g-1}/(g-1)!$ is algebraic on every principally
polarized abelian variety, so it stays of Hodge type everywhere. Yet by the
criterion of Matsusaka and Ran \cite{Mat59,Ran81} a curve with this class
exists only on Jacobians and on products of Jacobians, which form a proper
subvariety of the moduli space once $g\ge4$. By Bloch's theorem
\cite{Blo72} the Abel--Jacobi curve is therefore not semiregular. Schoen's
subvariety is the analogue one step up, built from the fibre of the
Abel--Jacobi map over the canonical class, and nothing in its construction
moves it off the locus of Prym varieties. This is an analogy and not a proof.
If the answer to Question~\ref{q:schoen} is no, the deformation route needs
a different object, and Theorem~\ref{thm:locus}(e) says it cannot come from
divisors, homomorphisms or the tautological cycles of curves.

The theorem also sits consistently against the barrier of \cite{BhW}.
Since $3-(1+\sigma+\sigma^{2})=(1-\sigma)(2+\sigma)$, the map $\Psi$ is
$(c_{i})\mapsto(1-\sigma)\sum_{i}(2+\sigma)(c_{i}-c_{0})$, which has the form
of the maps in \cite[Proposition~11]{BhW}; and $(1-\sigma)(J)=B$, because
$(2+\sigma)(2+\sigma^{2})=3$ on $(1-\sigma)(J)$, so that $2+\sigma$ is an
isogeny of it. That proposition says that $\Psi_{*}$
carries every class in the subring of $H^{*}(C^{h},\QQ)$ generated by the
diagonals, the graphs of $\sigma$ and $\sigma^{2}$ and the point classes to a
class invariant under the unitary group, and such a class pairs to zero with
the Weil plane. By Theorem~\ref{thm:schoensub}(a), then, $[\Lambda_{0}]$ does
not lie in that subring. This is what the barrier predicts: a Weil class at
a Prym variety with full Hodge group can come only from a cycle outside the
tautological ring, and Schoen's splitting of the canonical system into three
components produces one.

For fields other than $\QQ(\sqrt{-3})$ and $\QQ(\sqrt{-1})$, and for
nonsplit forms when $2n\ge6$, no cycle is known at any member with full
Hodge group, so the route has nothing to deform.

\section{Classes beyond the Weil classes}\label{sec:beyondweil}

By Theorem~\ref{thm:shape}(b), the conjecture for abelian varieties also
needs the classes that divisors and Weil classes do not generate. Their
first instance is Mumford's.

Let $A$ be a Mumford fourfold, so that $\End(A)=\ZZ$ and $\Hg(A)$ is a
$\QQ$-form of $\mathrm{SL}_{2}^{3}$ acting on $V_{\CC}=S_{1}\otimes S_{2}\otimes
S_{3}$. These fourfolds form families over Shimura curves
\cite{Mum69}. The space $\bigwedge^{2}V_{\CC}$ contains
$\mathrm{Sym}^{2}S_{1}\otimes\mathrm{Sym}^{2}S_{2}\otimes\mathbf{1}$ and its two
companions, and on $A\times A$ the invariants of $\Hg(A)$ include, besides the
classes generated by divisors, the two classes that come from the identity of
each of these three pieces modulo the polarization. They are Hodge classes
that are neither Lefschetz nor Weil classes. For them \cite{BhH8} proves:
they are algebraic at every CM point of the family; they are algebraic on
every member once one Kuga--Satake class, on a K3 surface of Picard number
thirteen attached to the fourfold, is algebraic; and they are algebraic on
every member exactly when the Lefschetz standard conjecture holds for one
ninefold. It also shows that no construction from line bundles by cones,
tensor products, pull-backs, push-forwards and Fourier--Mukai functors
reaches them, because all of these preserve invariance under the Lefschetz
group.

Here the attempt stops at one class. The Kuga--Satake correspondence is
known to be algebraic for a few families of K3 surfaces (see
Section~\ref{sec:nonabelian}), and the K3 surfaces attached to Mumford's
fourfolds are not among them.

\section{Varieties that are not of abelian type}\label{sec:nonabelian}

The statement (F3$'$) holds trivially on every variety whose cohomology is
reached from abelian varieties by algebraic correspondences, and
Proposition~\ref{prop:descent} spreads it along surjective morphisms. The
literature adds a good deal.

\begin{itemize}
\item \emph{Dimension at most three, and codegrees $0$ and $2$} hold by
Corollary~\ref{cor:smalldim}.
\item \emph{K3 surfaces of Picard number $19$ or $20$} carry a Shioda--Inose
structure, which relates their transcendental lattice to that of an abelian
surface \cite{Mor84}.
\item \emph{Products of surfaces.} The conjecture holds for every product of
surfaces with $q=2$ and $p_{g}=1$, and for every self-product of a K3 surface
whose field of Hodge endomorphisms of the transcendental lattice is a CM field
\cite{RM08}. It holds for every power of a K3 surface associated with a
variety of generalized Kummer type \cite{Fl25}, and for the squares of the
K3 surfaces in an explicit family covered by squares of curves of genus
seven \cite{ILP21}.
\item \emph{Hyper-K\"ahler varieties.} Rational Hodge isometries between
varieties of $K3^{[n]}$ type are algebraic \cite{Mar22}, and the standard
conjectures hold for them \cite{CM13}. The Chow motive of a generalized Kummer
variety lies in the tensor category generated by the abelian surface, so the
conjecture holds for all products of generalized Kummer varieties
\cite{Xu15,FTV19}. The same holds for O'Grady's six-dimensional varieties
\cite{Fl23}, and O'Grady's ten-dimensional varieties have motives built from
those of K3 surfaces \cite{FFZ21}. The motive of a smooth moduli space of
stable sheaves on a K3 or abelian surface $S$ is a summand of a sum of twisted
motives of powers of $S$ \cite{Bul18}.
\item \emph{Fourfolds.} The conjecture holds for uniruled fourfolds
\cite{CM78} and for cubic fourfolds \cite{Zuc77}, and cubic fourfolds whose
lattice of algebraic $2$-cycles has rank greater than $19$ have abelian
motive \cite{ABLP20}. Every smooth Delsarte hypersurface, and every cyclic
cover branched along one, satisfies (F3$'$) \cite{BhH8}.
\end{itemize}

The first variety on which (F3$'$) is open is the square $S\times S$ of a K3
surface $S$ whose transcendental lattice has real multiplication by a totally
real field $E\ne\QQ$ \cite{BhH8}. The endomorphisms in $E$ give Hodge classes
on $S\times S$ in $H^{2}(S)\otimes H^{2}(S)$ that are not generated by divisors.
If the Kuga--Satake correspondence between $S$ and its Kuga--Satake abelian
variety were algebraic, these classes would move to Hodge classes on the
square of an abelian variety, and (F3$'$) for $S\times S$ would become a case
of (F2). It is not known to be algebraic for such $S$. The analysis of
\cite{BhH8} gives the same structure here as for Weil classes: the algebraic
locus of the real multiplication in a family of such K3 surfaces is dense,
contains every CM point, and is either the whole family or meagre. Deciding
between the two is open.

Beyond the K3 surfaces lies the general case of (F3$'$): a variety of
dimension at least four whose motive is not known to come from abelian
varieties, with Hodge classes in a degree that Corollary~\ref{cor:smalldim}
does not reach. The methods that have produced cycles in the known cases are
few: divisors, through Theorem~\ref{thm:lef11}; rational curves, as in
\cite{CM78}; normal functions over Lefschetz pencils, as for cubic fourfolds
\cite{Zuc77}; the inductive structure of Fermat and Delsarte varieties; and
correspondences with curves, abelian varieties and moduli of sheaves, as in
the results just listed. We do not see how to apply any of them to a
variety with none of this structure, and that is where the attempt at (F3$'$)
stops.

\section{Spreading Schoen's construction}\label{sec:spread}

Schoen's cycles reach a family of dimension $3n$ inside one of dimension
$n^{2}$, and Question~\ref{q:schoen} asks whether they can be deformed out of
it. This section tries the two other ways of spreading them: carrying them to
other abelian varieties by correspondences, and running the construction on
other varieties. The first gives a theorem. It makes the Weil class of any
abelian variety of Weil type for $\QQ(\sqrt{-3})$ algebraic exactly when the
Weil class of a second, complementary variety is, but a count of dimensions
shows that it cannot by itself reach the general member. The second reaches
neither (F2) nor (F3$'$), for reasons we make precise.

Throughout the section $K=\QQ(\sqrt{-3})$ and $\zeta$ is a primitive cube
root of unity. Varieties of Weil type are for $K$, with any discriminant,
and we assume that $\cO_{K}$ acts, which can always be arranged by passing
to a $K$-isogenous variety. For such a variety $A$ of dimension $2m$ write
$W(A)\subseteq H^{2m}(A,\QQ)$ for its Weil plane. Over $\CC$ the space
$H^{1}(A,\CC)$ is the sum of the eigenspaces $V_{\chi}(A)$ and
$V_{\bar\chi}(A)$ of $\zeta^{*}$, each of dimension $2m$, and
$W(A)\otimes\CC=\bigwedge^{2m}V_{\chi}(A)\oplus\bigwedge^{2m}V_{\bar\chi}(A)$.
We call $A$ \emph{Weil-algebraic} if every class in $W(A)$ is algebraic. One
nonzero algebraic class in $W(A)$ is enough, because a suitable element of
$\cO_{K}$ acting as an endomorphism moves it to a second class that spans
$W(A)$ with it \cite[proof of Theorem~6]{BhW}.

\subsection{Products and complements}

\begin{lemma}\label{lem:twothree}
Let $B$, $A_{1}$ and $A_{2}$ be of Weil type, of dimensions
$2(m_{1}+m_{2})$, $2m_{1}$ and $2m_{2}$, and let $f\colon B\to A_{1}\times
A_{2}$ be an isogeny that commutes with $\cO_{K}$.
\begin{enumerate}[label=(\alph*)]
\item The isomorphism $f^{*}$ carries a plane in
$W(A_{1})\otimes W(A_{2})\subseteq H^{2m_{1}}(A_{1})\otimes H^{2m_{2}}(A_{2})$
onto $W(B)$. Over $\CC$ this plane is spanned by $\omega_{1}\otimes\omega_{2}$
and $\bar\omega_{1}\otimes\bar\omega_{2}$, where $\omega_{i}$ spans
$\bigwedge^{2m_{i}}V_{\chi}(A_{i})$ and $\bar\omega_{i}$ is its complex
conjugate.
\item If two of $B$, $A_{1}$ and $A_{2}$ are Weil-algebraic, so is the third.
\end{enumerate}
\end{lemma}

\begin{proof}
The map $f^{*}$ is an isomorphism on rational cohomology that commutes with
$\cO_{K}$ and carries algebraic classes to algebraic classes, and so does its
inverse $(\deg f)^{-1}f_{*}$. So we may take $B=A_{1}\times A_{2}$ with the
diagonal action. Then $V_{\chi}(B)=V_{\chi}(A_{1})\oplus V_{\chi}(A_{2})$, and
the top exterior power of a direct sum is the tensor product of the top
exterior powers. This is~(a).

If $A_{1}$ and $A_{2}$ are Weil-algebraic, every class in
$W(A_{1})\otimes W(A_{2})$ is a sum of products
$\mathrm{pr}_{1}^{*}\alpha_{1}\cup\mathrm{pr}_{2}^{*}\alpha_{2}$ of algebraic
classes, so $B$ is Weil-algebraic by~(a). Suppose now that $B$ and $A_{2}$
are Weil-algebraic. On $W(A_{2})$ we have $\int\omega_{2}\cup\omega_{2}=0$,
because $\bigwedge^{4m_{2}}V_{\chi}(A_{2})=0$, while
$p:=\int\omega_{2}\cup\bar\omega_{2}\ne0$, because
$\bigwedge^{2m_{2}}V_{\chi}(A_{2})\otimes\bigwedge^{2m_{2}}V_{\bar\chi}(A_{2})$
is the top exterior power of $H^{1}(A_{2},\CC)$. The number $p$ is real,
since $\omega_{2}$ and $\bar\omega_{2}$ have even degree and so commute. Take nonzero rational
classes $x\in W(B)$ and $\gamma\in W(A_{2})$. Being real, they have the form
$x=s\,\omega_{1}\otimes\omega_{2}+\bar s\,\bar\omega_{1}\otimes\bar\omega_{2}$
and $\gamma=t\,\omega_{2}+\bar t\,\bar\omega_{2}$ with $s,t\ne0$. The class
\[
  \mathrm{pr}_{1*}\bigl(x\cup\mathrm{pr}_{2}^{*}\gamma\bigr)
  =p\,\bigl(s\bar t\;\omega_{1}+\bar s t\;\bar\omega_{1}\bigr)
\]
is rational, algebraic and nonzero, and it lies in $W(A_{1})$. So $A_{1}$ is
Weil-algebraic.
\end{proof}

\begin{theorem}\label{thm:transfer}
Let $A$ be an abelian variety of Weil type for $K$, of dimension $2m\ge2$,
on which $\cO_{K}$ acts. Then there are a smooth curve $C\subseteq A$ stable
under $\zeta$, on which $\zeta$ acts freely, so that $\pi\colon C\to X=C/\zeta$
is an \'etale cyclic triple cover of a curve of some genus $g\ge m+1$; a
variety $A'$ of Weil type of dimension $2(g-1-m)$; and an isogeny commuting
with $\cO_{K}$ between the Prym variety $B$ of $\pi$ and $A\times A'$. For
every such curve, $A$ is Weil-algebraic if and only if $A'$ is.
\end{theorem}

\begin{proof}
\emph{The curve.} The automorphism $\zeta$ of $A$ has finitely many fixed
points, the kernel of the isogeny $1-\zeta$. Let $L$ be ample on $A$. The
bundle $N=L\otimes\zeta^{*}L\otimes\zeta^{2*}L$ satisfies $\zeta^{*}N\cong N$,
so, the group being cyclic, it carries a $\zeta$-linearization, and the
linearized bundle $N^{3}$ descends to an ample line bundle $\bar M$ on the
normal projective quotient $q\colon A\to A/\zeta$, whose singular points are
among the images of the fixed points. For $k$ large, $2m-1$ general members of $|\bar M^{k}|$ meet in a
smooth irreducible curve $X$ that avoids those finitely many points, by
Bertini's theorem. Then $C=q^{-1}(X)$ is \'etale over $X$, hence smooth, and
it is the complete intersection of the $2m-1$ ample divisors $q^{*}D_{i}$,
hence connected. So $\pi\colon C\to X$ is an \'etale cyclic triple cover, and
$\zeta$ acts on $C$ as its covering automorphism $\sigma$. Schoen observes
the existence of such curves in \cite[Section~3]{Sch88}.

\emph{The surjection.} By the Lefschetz hyperplane theorem, restriction
$H^{1}(A,\QQ)\to H^{1}(C,\QQ)$ is injective. So the homomorphism
$\alpha\colon J(C)\to A$ that adds up the points of a divisor of degree zero
is surjective, and it commutes with $\zeta$. In $\cO_{K}$ we have
$1+\zeta+\zeta^{2}=0$, so $\alpha\circ(1+\sigma+\sigma^{2})=0$ and
$\alpha\circ P=3\alpha$ for $P=3-(1+\sigma+\sigma^{2})$. Hence $\alpha$ maps
the Prym variety $B=P(J(C))$ onto $A$, and in particular $g-1\ge m$.

\emph{The complement.} Let $A'$ be the connected component of the kernel of
$\alpha|_{B}$, and $A''$ its complement in $B$ for the polarization induced
from $J(C)$. Since $\sigma$ preserves that polarization, both are stable under
$\sigma$, so $\cO_{K}$ acts on them. The maps $A''\to A$ and
$A'\times A''\to B$ are isogenies commuting with $\cO_{K}$, and composing one
with a quasi-inverse of the other relates $B$ and $A\times A'$.

\emph{Weil type.} By \cite[Lemma~1.6a]{Sch88} the two eigenspaces of $\sigma$
on $H^{1,0}(B)$ both have dimension $g-1$, and on $H^{1,0}(A)$ both have
dimension $m$. So on $H^{1,0}(A')$ both have dimension $g-1-m$, and $A'$ is
of Weil type.

\emph{The equivalence.} The cover is \'etale, so Schoen's Theorem~2.0 applies
to it with $r=0$, and by \cite[Corollary~3.1]{Sch88} the variety $B$ is
Weil-algebraic. Lemma~\ref{lem:twothree}(b) gives the equivalence.
\end{proof}

When $g-1=m$ the complement is zero and $A$ is itself $K$-isogenous to a
Prym variety; that happens on the Prym locus, of dimension $3m$, and nowhere
else. For larger $g$ the theorem trades the Weil class of $A$ for that of
$A'$, and since curves of higher degree give complements of any size, one
can hope to choose the curve so that $A'$ is a variety on which the Weil
class is already known.

\begin{corollary}\label{cor:transfer}
$A$ is Weil-algebraic as soon as, for one curve as in
Theorem~\ref{thm:transfer}, the complement $A'$ is $K$-isogenous to a product
of Weil-algebraic varieties. Among these are Prym varieties of \'etale cyclic
triple covers; abelian surfaces of Weil type; Weil fourfolds for $K$
\cite{Sch88,Sch98}; and the varieties $E^{k}\times\bar E^{k}$, where $E$ is an
elliptic curve with complex multiplication by $\cO_{K}$ and $\bar E$ is $E$
with $\cO_{K}$ acting through complex conjugation.
\end{corollary}

\begin{proof}
Products of Weil-algebraic varieties are Weil-algebraic by
Lemma~\ref{lem:twothree}(b), and Theorem~\ref{thm:transfer} finishes the
argument. On an abelian surface the Weil plane lies in $H^{2}$ and consists of
Hodge classes, which are algebraic by Theorem~\ref{thm:lef11}. On
$E^{k}\times\bar E^{k}$, with coordinates $z_{1},\dots,z_{2k}$ on the factors,
$V_{\chi}$ is spanned by $dz_{1},\dots,dz_{k}$ and
$d\bar z_{k+1},\dots,d\bar z_{2k}$, so up to sign
\[
  \omega_{\chi}=\bigwedge_{i=1}^{k}\bigl(dz_{i}\wedge d\bar z_{k+i}\bigr).
\]
Each factor is a class of type $(1,1)$ on $E\times\bar E$, whose N\'eron--Severi
group has rank $4=h^{1,1}$ because $E$ has complex multiplication, so it is a
complex combination of divisor classes. Hence $\omega_{\chi}$ and its
conjugate are polynomials in divisor classes.
\end{proof}

\begin{remark}\label{rem:count}
For a general $A$ with $m\ge4$ one cannot expect a complement of these
kinds. Pryms of \'etale triple covers of curves of genus $m+k+1$ form a
family of dimension $3(m+k)$ inside the Weil family of dimension $(m+k)^{2}$.
The varieties $K$-isogenous to $A\times A'$, with $A$ in the family of
dimension $m^{2}$ and $A'$ in a family of dimension $f$, form a subvariety of
dimension $m^{2}+f$. For $A$ to be general, the two must meet in dimension at
least $m^{2}$, which needs
\[
  3(m+k)+f\ \ge\ (m+k)^{2}
\]
unless they meet in more than the expected dimension. Products of the
varieties in Corollary~\ref{cor:transfer} of total dimension $2k$ have at most
$3k$ moduli, and $3m+6k<(m+k)^{2}$ for every $k\ge0$ once $m\ge4$. For
$m=3$ only $k=0$ passes, which is the Prym locus itself, and for $m=2$ the
values $k\le2$ pass. So, unless the Prym locus meets these product loci in
excess dimension, the corollary reaches only a meagre subset of the family
once $m\ge4$, as Schoen anticipated \cite[Section~3]{Sch88}. Such excess
intersections are what the Zilber--Pink philosophy says should not happen.
This is a count and not a proof.
\end{remark}

\subsection{Other varieties, and \texorpdfstring{(F2) and (F3$'$)}{(F2) and (F3')}}

\emph{(F2).} Schoen's classes are Weil classes. By \cite[Lemma~1.2 and
Remark~1.4]{Sch88}, for every curve with a cyclic group of automorphisms the
space of classes his Theorem~2.0 makes algebraic is the pull-back of the Weil
plane of a generalized Prym variety with cyclotomic multiplication. For
cyclotomic fields of degree at least four this bears on the third missing
input of \cite{BhH8}, the propagation for CM fields of higher degree. It
does not bear on (F2). Schoen's own count in \cite[Section~3]{Sch88} shows
that these generalized Pryms can fill a whole Weil family in at most four
small cases. To reach a class of (F2), such as the exceptional classes on $A\times A$
for a Mumford fourfold $A$, the construction would need an algebraic
correspondence from a Prym-type variety $B$ to $A\times A$ that carries a Weil
class of $B$ to the exceptional class. That correspondence is a Hodge class on
$B\times A\times A$, so asking for it is again an instance of the conjecture
for abelian varieties. Nor can the construction be run on $A$ itself: a
Mumford fourfold has $\End(A)=\ZZ$, so it has no automorphism of order
three.

\emph{(F3$'$).} Schoen defines the same space of classes for any smooth
projective $C$ of odd dimension $k$ with a cyclic group of automorphisms,
under a few hypotheses, the main one being that the part $V$ of
$H^{k}(C,\QQ)$ on which the group acts by primitive characters has only two
Hodge types. His example is a quintic threefold
invariant under $\ZZ/5$ and missing the fixed points, where the space lies in
$H^{120}(C^{40},\QQ)$, has type $(60,60)$, and was not known to be algebraic
\cite[Section~1]{Sch88}. In that example $H^{3,0}$ is invariant, so $V$ has
types $(2,1)$ and $(1,2)$. It is then of level one, and so it is a Tate
twist of $H^{1}$ of an abelian variety $J_{V}$, and the classes are images of
the Weil plane of $J_{V}$ under a Hodge isomorphism. They would be algebraic
if $J_{V}$ were Weil-algebraic and the isomorphism were induced by an
algebraic correspondence. The second requirement is the generalized Hodge
conjecture for the level-one part of $H^{3}(C)$, an instance of (F3$'$) that
we do not know to be proved for these quintics. The cycles of \cite[Section~2]{Sch88} come from
the Abel--Jacobi map of a curve, which supplies that correspondence for free.
For a variety of higher dimension nothing supplies it, so the construction
stops at curves. When the two Hodge types are not adjacent, $V$ is not of
level one, and the construction has no abelian variety to work on.

So spreading Schoen's construction by correspondences reduces the Weil class
of $A$ to that of a complement without making it easier, and spreading it to
other varieties leads back to (F2) and (F3$'$) themselves. Neither closes
anything that Question~\ref{q:schoen} leaves open.

\section{What is proved and what is not}\label{sec:end}

Proved here, with complete proofs or by citation of refereed work:
\begin{enumerate}[label=(\arabic*)]
\item every Hodge class of degree $2$ or $2n-2$, and every Hodge class on a
variety of dimension at most three, is algebraic;
\item the conjecture descends along surjective morphisms, and reduces to
middle degree on complete intersections given the Lefschetz standard
conjecture;
\item on every Prym variety $B$ of an \'etale cyclic triple cover of a
curve of genus $n+1$, Schoen's construction gives an irreducible
$n$-dimensional subvariety $Y$ whose class pairs nontrivially with the Weil
plane, and at a very general such $B$ its class is $c\eta^{n}+w$ with $c>0$
and $w$ a nonzero Weil class (Theorem~\ref{thm:schoensub});
\item if that subvariety is semiregular for one very general cover of genus
$n+1$, the Weil classes are algebraic on every member of the split Weil
family of $\QQ(\sqrt{-3})$ in dimension $2n$ (Corollary~\ref{cor:schoensemireg});
\item every abelian variety $A$ of Weil type for $\QQ(\sqrt{-3})$ is, up to
isogeny, a factor $A\times A'$ of the Prym variety of an \'etale cyclic
triple cover, with $A'$ of Weil type, and $A$ is Weil-algebraic if and only
if $A'$ is (Theorem~\ref{thm:transfer}); in particular $A$ is Weil-algebraic
whenever $A'$ is a product of the varieties listed in
Corollary~\ref{cor:transfer}.
\end{enumerate}

Recalled from \cite{BhH8} and used as the frame: the conjecture is equivalent
to (F2) and (F3$'$) together; (F2) is the conjecture for abelian varieties;
and the minimal sets of hypotheses are those of Theorem~\ref{thm:shape}(c).

Not proved here, and open:
\begin{enumerate}[label=(\roman*)]
\item the Lefschetz standard conjecture, in general;
\item the Weil classes on a general abelian variety of Weil type of
dimension at least eight, and of dimension six outside the known families;
in particular Question~\ref{q:schoen}, and the existence, for a general $A$
of dimension at least eight, of a complement as in
Theorem~\ref{thm:transfer} on which the Weil classes are known, which the
count of Remark~\ref{rem:count} makes unlikely;
\item the exceptional classes on the square of a Mumford fourfold, which
follow from one Kuga--Satake class;
\item the conjecture for the square of a K3 surface with real
multiplication, the first open case of (F3$'$), and (F3$'$) in general.
\end{enumerate}

So the attempt in full ends where it was always likely to end. The
conjecture is equivalent to two statements, and each of them has an explicit
first open case. The most concrete thing this paper adds is that the Weil
case for $\QQ(\sqrt{-3})$ in every dimension now hangs on a single, explicitly
described subvariety of a Prym variety and the question whether it deforms.

\subsection*{Machine checks}
No proof in this paper depends on a computer. The finite parts of the
arguments have nevertheless been repeated by machine, and the files are in
the author's repository \cite{BhHC}.

In Lean~4, checked by its kernel with nothing beyond the core library, we
verified five things. The first is the deductions of
Propositions~\ref{prop:hardlef}, \ref{prop:descent} and~\ref{prop:middle} and
of Corollary~\ref{cor:smalldim}, with every geometric input stated as a
hypothesis, so that the dependence of Proposition~\ref{prop:middle} on the
Lefschetz standard conjecture is visible in the statement. The second is the
counts of genera and dimensions in Section~\ref{sec:weil}, for every genus at
once. The third is the arithmetic of the characters of $\ZZ/3$ in $\ZZ[\zeta]$
used in Step~3 of Theorem~\ref{thm:schoensub}. The fourth is the weight count
of Step~5 and the deduction of $c>0$ and $w\ne0$ from it. The identity
$3-(1+\sigma+\sigma^{2})=(1-\sigma)(2+\sigma)$ used after
Question~\ref{q:schoen} is checked there too. The fifth, for
Section~\ref{sec:spread}, is the dimensions in Theorem~\ref{thm:transfer}, including the
equality of the two eigenspaces on $A'$, the relation $\alpha\circ P=3\alpha$,
and the arithmetic of Remark~\ref{rem:count} for every $m$ and $k$ at once.
It checks only that arithmetic; the remark itself remains a count and not a
proof.

In Julia, and in a Python version of the same checks, all in exact
arithmetic, we computed the following. On a general abelian variety of Weil
type of dimension $2n\le8$, the Hodge classes span one line in each even
degree except the middle one, where they span three. For a Weil class $w$
the products $\eta\cup w$ and $\eta^{n}\cup w$ vanish. The space $U_{\chi}$
of Step~1 is one-dimensional for $g=3$ and $g=4$; for $g=3$ this was checked
by summing over the whole group. For $h\le6$ the class $u_{\chi}$ has $h!$
terms, is symmetric, and pairs nontrivially with $u_{\bar\chi}$. For
Lemma~\ref{lem:twothree} with $2m_{1},2m_{2}\in\{2,4\}$ we computed the
push-forward of $x\cup\mathrm{pr}_{2}^{*}\gamma$ in the exterior algebra and
found the formula of the proof, and for Corollary~\ref{cor:transfer} with
$k\le4$ we checked that $\omega_{\chi}$ is, up to sign, the product of the
classes $dz_{i}\wedge d\bar z_{k+i}$.

None of this touches the geometry the paper relies on: the theorems of
Schoen, Bloch, Looijenga and Andr\'e, and those of \cite{BhH8}. None of it
bears on Question~\ref{q:schoen}.

\subsection*{A note on preparation}
Drafts of parts of this paper were prepared with the help of an AI language
model, used under the author's direction. The author has checked every
statement and proof and is responsible for the content.

\begin{thebibliography}{ABLP20}

\bibitem[ABLP20]{ABLP20}
H.~Awada, M.~Bolognesi, R.~Laterveer and C.~Pedrini, \emph{K3 surfaces and
cubic fourfolds with abelian motive}, preprint, arXiv:2012.05969.

\bibitem[AH62]{AH62}
M.~F.~Atiyah and F.~Hirzebruch, \emph{Analytic cycles on complex manifolds},
Topology \textbf{1} (1962), 25--45.

\bibitem[And92]{And92}
Y.~Andr\'e, \emph{Mumford--Tate groups of mixed Hodge structures and the
theorem of the fixed part}, Compositio Math.\ \textbf{82} (1992), 1--24.

\bibitem[And96]{And96}
Y.~Andr\'e, \emph{Pour une th\'eorie inconditionnelle des motifs}, Inst.\
Hautes \'Etudes Sci.\ Publ.\ Math.\ \textbf{83} (1996), 5--49.

\bibitem[BhH8]{BhH8}
D.~Bhattacharjee, \emph{Algebraic Loci of Weil Classes on Abelian
Varieties}, preprint, version 5.1.0, Zenodo, 2026,
doi:10.5281/zenodo.23227940.

\bibitem[BhHC]{BhHC}
D.~Bhattacharjee, \emph{Hodge-Conjecture-Full}, repository of the Lean and
Julia checks for this paper, 2026,
\url{https://github.com/creelie/Hodge-Conjecture-Full}, directory
\texttt{verification}.

\bibitem[BhW]{BhW}
D.~Bhattacharjee, \emph{Closing the algebraic locus of Weil classes: an
attempt and the exact point where it stops}, preprint, 2026.

\bibitem[Blo72]{Blo72}
S.~Bloch, \emph{Semi-regularity and de Rham cohomology}, Invent.\ Math.\
\textbf{17} (1972), 51--66.

\bibitem[B\"ul18]{Bul18}
T.-H.~B\"ulles, \emph{Motives of moduli spaces on K3 surfaces and of special
cubic fourfolds}, preprint, arXiv:1806.08284.

\bibitem[CM13]{CM13}
F.~Charles and E.~Markman, \emph{The standard conjectures for holomorphic
symplectic varieties deformation equivalent to Hilbert schemes of K3
surfaces}, preprint, arXiv:1009.0413.

\bibitem[CM78]{CM78}
A.~Conte and J.~P.~Murre, \emph{The Hodge conjecture for fourfolds admitting
a covering by rational curves}, Math.\ Ann.\ \textbf{238} (1978), 79--88.

\bibitem[Del82]{Del82}
P.~Deligne, \emph{Hodge cycles on abelian varieties} (notes by J.~S.~Milne),
in: Hodge Cycles, Motives, and Shimura Varieties, Lecture Notes in Math.\
\textbf{900}, Springer, Berlin, 1982, 9--100.

\bibitem[FFZ21]{FFZ21}
S.~Floccari, L.~Fu and Z.~Zhang, \emph{On the motive of O'Grady's
ten-dimensional hyper-K\"ahler varieties}, Commun.\ Contemp.\ Math.\
\textbf{23} (2021), 2050034.

\bibitem[Fl23]{Fl23}
S.~Floccari, \emph{On the motive of O'Grady's six-dimensional hyper-K\"ahler
varieties}, \'Epijournal G\'eom.\ Alg\'ebrique \textbf{7} (2023), Article~4.

\bibitem[Fl25]{Fl25}
S.~Floccari, \emph{K3 surfaces associated with varieties of generalized
Kummer type}, preprint, arXiv:2501.02315.

\bibitem[FTV19]{FTV19}
L.~Fu, Z.~Tian and C.~Vial, \emph{Motivic hyper-K\"ahler resolution
conjecture I: generalized Kummer varieties}, preprint, arXiv:1608.04968.

\bibitem[ILP21]{ILP21}
C.~Ingalls, A.~Logan and O.~Patashnick, \emph{Explicit coverings of families
of elliptic surfaces by squares of curves}, preprint, arXiv:2009.07807.

\bibitem[Kle68]{Kle68}
S.~L.~Kleiman, \emph{Algebraic cycles and the Weil conjectures}, in: Dix
Expos\'es sur la Cohomologie des Sch\'emas, North-Holland, Amsterdam, 1968,
359--386.

\bibitem[Koi04]{Koi04}
K.~Koike, \emph{Algebraicity of some Weil Hodge classes}, Canad.\ Math.\
Bull.\ \textbf{47} (2004), 566--572.

\bibitem[Kol92]{Kol92}
J.~Koll\'ar, in: E.~Ballico, F.~Catanese and C.~Ciliberto (eds.),
\emph{Trento examples}, in: Classification of Irregular Varieties, Lecture
Notes in Math.\ \textbf{1515}, Springer, Berlin, 1992, 134--135.

\bibitem[KS53]{KS53}
K.~Kodaira and D.~C.~Spencer, \emph{Divisor class groups on algebraic
varieties}, Proc.\ Nat.\ Acad.\ Sci.\ U.S.A.\ \textbf{39} (1953), 872--877.

\bibitem[Lef24]{Lef24}
S.~Lefschetz, \emph{L'analysis situs et la g\'eom\'etrie alg\'ebrique},
Gauthier-Villars, Paris, 1924.

\bibitem[Lie68]{Lie68}
D.~I.~Lieberman, \emph{Numerical and homological equivalence of algebraic
cycles on Hodge manifolds}, Amer.\ J.\ Math.\ \textbf{90} (1968), 366--374.

\bibitem[Loo97]{Loo97}
E.~Looijenga, \emph{Prym representations of mapping class groups}, Geom.\
Dedicata \textbf{64} (1997), 69--83.

\bibitem[Mar22]{Mar22}
E.~Markman, \emph{Rational Hodge isometries of hyper-K\"ahler varieties of
$K3^{[n]}$-type are algebraic}, preprint, arXiv:2204.00516.

\bibitem[Mar25a]{Mar25a}
E.~Markman, \emph{Cycles on abelian $2n$-folds of Weil type from secant
sheaves on abelian $n$-folds}, preprint, arXiv:2502.03415.

\bibitem[Mat59]{Mat59}
T.~Matsusaka, \emph{On a characterization of a Jacobian variety}, Mem.\
Coll.\ Sci.\ Univ.\ Kyoto Ser.\ A Math.\ \textbf{32} (1959), 1--19.

\bibitem[Mor84]{Mor84}
D.~R.~Morrison, \emph{On K3 surfaces with large Picard number}, Invent.\
Math.\ \textbf{75} (1984), 105--121.

\bibitem[Mum69]{Mum69}
D.~Mumford, \emph{A note of Shimura's paper ``Discontinuous groups and
abelian varieties''}, Math.\ Ann.\ \textbf{181} (1969), 345--351.

\bibitem[MZ99]{MZ99}
B.~Moonen and Yu.~G.~Zarhin, \emph{Hodge classes on abelian varieties of low
dimension}, Math.\ Ann.\ \textbf{315} (1999), 711--733.

\bibitem[PZ25]{PZ25}
D.~Patel and Y.~Zhang, \emph{Algebraicity of Hodge classes on some
generalized Prym varieties}, preprint, arXiv:2506.13729.

\bibitem[Ran81]{Ran81}
Z.~Ran, \emph{On subvarieties of abelian varieties}, Invent.\ Math.\
\textbf{62} (1981), 459--479.

\bibitem[RM08]{RM08}
J.~J.~Ram\'on Mar\'i, \emph{On the Hodge conjecture for products of certain
surfaces}, Collect.\ Math.\ \textbf{59} (2008), 1--26.

\bibitem[Sch88]{Sch88}
C.~Schoen, \emph{Hodge classes on self-products of a variety with an
automorphism}, Compositio Math.\ \textbf{65} (1988), 3--32.

\bibitem[Sch98]{Sch98}
C.~Schoen, \emph{Addendum to: Hodge classes on self-products of a variety
with an automorphism}, Compositio Math.\ \textbf{114} (1998), 329--336.

\bibitem[Voi02b]{Voi02b}
C.~Voisin, \emph{A counterexample to the Hodge conjecture extended to
K\"ahler varieties}, Int.\ Math.\ Res.\ Not.\ \textbf{2002}, no.~20,
1057--1075.

\bibitem[Wei77]{Wei77}
A.~Weil, \emph{Abelian varieties and the Hodge ring}, in: \OE uvres
Scientifiques, Collected Papers, Vol.~III, Springer, New York, 1979,
421--429.

\bibitem[Xu15]{Xu15}
Z.~Xu, \emph{Algebraic cycles on a generalized Kummer variety}, preprint,
arXiv:1506.04297.

\bibitem[Zuc77]{Zuc77}
S.~Zucker, \emph{The Hodge conjecture for cubic fourfolds}, Compositio
Math.\ \textbf{34} (1977), 199--209.

\end{thebibliography}

\end{document}
